\chapter{DX-034 Declarative View Data Binding}\label{record-132}

\noindent{\footnotesize\color{BookMuted}\path{docs/design/DX-034-declarative-view-data-binding.md}\par}\smallskip

\begin{RecordMeta}
title: DX-034 Declarative View Data Binding
legend: DX
lane: design

\end{RecordMeta}

\emph{Cycle for formalizing unidirectional, declarative, immutable data binding between business logic, active views, and block composition}

Legend:

\begin{itemize}
\tightlist
\item
  \href{https://github.com/flyingrobots/bijou/blob/1853ad902f6bf219326409e7f743a5a1d6517dbc/docs/legends/DX-developer-experience.md}{DX - Developer Experience}
\end{itemize}

Depends on:

\begin{itemize}
\tightlist
\item
  \href{https://github.com/flyingrobots/bijou/blob/cbd70612140e7a450d764b3c7be1c6af878b99bc/docs/design/DX-031-standard-bijou-blocks.md}{DX-031 - Standard Bijou Blocks}
\item
  \href{https://github.com/flyingrobots/bijou/blob/cbd70612140e7a450d764b3c7be1c6af878b99bc/docs/design/RE-001-define-runtime-engine-architecture.md}{RE-001 - Define Runtime Engine Architecture}
\item
  \href{https://github.com/flyingrobots/bijou/blob/cbd70612140e7a450d764b3c7be1c6af878b99bc/docs/design/RE-003-retain-layout-trees-and-layout-invalidation.md}{RE-003 - Retain Layout Trees And Layout Invalidation}
\item
  \href{https://github.com/flyingrobots/bijou/blob/cbd70612140e7a450d764b3c7be1c6af878b99bc/docs/design/RE-005-buffer-commands-and-effects-separately.md}{RE-005 - Buffer Commands And Effects Separately}
\item
  \href{https://github.com/flyingrobots/bijou/blob/cbd70612140e7a450d764b3c7be1c6af878b99bc/docs/design/RE-006-formalize-component-layout-and-interaction-contracts.md}{RE-006 - Formalize Component Layout And Interaction Contracts}
\end{itemize}

\section{Why This Cycle Exists}\label{record-132-why-this-cycle-exists}

DX-031 gave Bijou a metadata-first block contract, but it did not answer the next developer-experience question:

How does a block or view declare the data it needs, and how does that data arrive when the block is part of the active view hierarchy?

Bijou already has adjacent runtime pieces:

\begin{itemize}
\tightlist
\item
  explicit application state and view stacks
\item
  retained layouts for active views
\item
  component layout and interaction contracts
\item
  command and effect buffers
\item
  an event bus for runtime messages
\end{itemize}

Those are necessary, but they are not a data-binding architecture.

This cycle exists to formalize the missing layer before AppShell, nested blocks, schema-bound blocks, or generated block packages harden the wrong API.

\section{Core Rule}\label{record-132-core-rule}

Data binding in Bijou is:

\begin{itemize}
\tightlist
\item
  unidirectional
\item
  declarative
\item
  immutable
\end{itemize}

Data flows from business logic and provider adapters into active views.

Views render immutable snapshots.

Views do not mutate business state, provider state, caches, stores, channels, or resources.

Views communicate user intent by emitting Commands. Business logic handles those Commands, updates state or provider inputs, and produces the next immutable snapshot. That is the loop.

% Presentation of original arrow notation. See flow-manifest.json.
\begin{center}
\begin{tikzpicture}[book flow node/.style={draw=BookHeading,line width=.65pt,fill=BookPanel,rounded corners=1.5pt,align=center,inner sep=5pt,minimum height=10mm,font=\fontsize{9}{11}\selectfont,execute at begin node={\hyphenpenalty=10000\exhyphenpenalty=10000}},book flow edge/.style={-{Latex[length=2mm]},draw=BookHeading,line width=.7pt}]
\node[book flow node,text width=96mm] (n0) at (0,0.0) {business state /\allowbreak{} adapters};
\node[book flow node,text width=96mm] (n1) at (0,-1.55) {providers};
\node[book flow node,text width=96mm] (n2) at (0,-3.1) {immutable binding snapshots};
\node[book flow node,text width=96mm] (n3) at (0,-4.65) {active view hierarchy};
\node[book flow node,text width=96mm] (n4) at (0,-6.2) {render};
\node[book flow node,text width=96mm] (n5) at (0,-7.75) {user input};
\node[book flow node,text width=96mm] (n6) at (0,-9.3) {Commands expressing intent};
\node[book flow node,text width=96mm] (n7) at (0,-10.85) {business logic};
\node[book flow node,text width=96mm] (n8) at (0,-12.4) {next state /\allowbreak{} provider inputs};
\draw[book flow edge] (n0) -- (n1);
\draw[book flow edge] (n1) -- (n2);
\draw[book flow edge] (n2) -- (n3);
\draw[book flow edge] (n3) -- (n4);
\draw[book flow edge] (n4) -- (n5);
\draw[book flow edge] (n5) -- (n6);
\draw[book flow edge] (n6) -- (n7);
\draw[book flow edge] (n7) -- (n8);
\end{tikzpicture}
\end{center}


There is no reverse data mutation path from view to provider.

\section{Boundary}\label{record-132-boundary}

This cycle is about data requirements and binding lifecycle.

It is not about:

\begin{itemize}
\tightlist
\item
  two-way binding
\item
  mutable stores exposed to views
\item
  hidden global provider registration
\item
  letting views read process, filesystem, network, or database adapters directly
\item
  making schemas own presentation
\item
  replacing Commands with arbitrary event side effects
\item
  forcing every component to become provider-aware
\end{itemize}

\section{Sponsor Human}\label{record-132-sponsor-human}

App builders composing shell, reader, inspector, form, and data blocks who need the view tree to bind to business data without manual prop plumbing or hidden mutation.

\section{Sponsor Agent}\label{record-132-sponsor-agent}

Agents generating Bijou apps or block packages that need to inspect what data a view requires, which providers can satisfy it, and which Commands the view can emit.

\section{Hill}\label{record-132-hill}

A builder can declare an AppShell whose navigation, content, inspector, and status regions bind to provider-backed immutable data snapshots, while every user action leaves the view tree as an explicit Command intent handled by business logic.

\section{Playback Questions}\label{record-132-playback-questions}

\begin{itemize}
\tightlist
\item
  Which data requirements does this active view tree declare?
\item
  Which provider satisfies each requirement?
\item
  What is the immutable snapshot version currently rendered?
\item
  Is each binding ready, loading, stale, empty, or failed?
\item
  Which Commands can this view emit?
\item
  Which business reducer or adapter action handles each Command?
\item
  What provider subscriptions are active only because this view is active?
\item
  What is disposed or suspended when the view leaves the active hierarchy?
\end{itemize}

\section{Vocabulary}\label{record-132-vocabulary}

\subsection{Business Logic}\label{record-132-business-logic}

The state machine, reducers, command handlers, domain services, and host adapters that own application truth.

Business logic may read or write external systems through ports and adapters. Views may not.

\subsection{Provider}\label{record-132-provider}

A runtime-backed source that can produce immutable snapshots for a declared data requirement.

A provider may be backed by:

\begin{itemize}
\tightlist
\item
  in-memory business state
\item
  a channel or message stream
\item
  a query cache
\item
  a filesystem adapter
\item
  a database adapter
\item
  a GraphQL, REST, or MCP adapter
\item
  a localization catalog
\item
  a credential-status adapter that never exposes secret values
\end{itemize}

The backing mechanism is not visible to views. The view sees only a binding snapshot.

\subsection{Data Requirement}\label{record-132-data-requirement}

A declarative statement made by a block, view, or component:

\begin{CodeBlock}
I need this named resource or signal, with these parameters, in this slot.
\end{CodeBlock}

Requirements are data, not callbacks. They are inspectable before binding.

\subsection{Binding Scope}\label{record-132-binding-scope}

The provider registry visible to a subtree of the active view hierarchy.

Scopes let an app provide one navigation source at the shell level, override a selection source inside a modal, or bind a test provider for a story without touching global state.

DX-034B makes scopes explicit local values in the public contract. A \texttt{P\allowbreak{}r\allowbreak{}o\allowbreak{}v\allowbreak{}i\allowbreak{}d\allowbreak{}e\allowbreak{}r\allowbreak{}S\allowbreak{}c\allowbreak{}o\allowbreak{}p\allowbreak{}e} declares which \texttt{D\allowbreak{}a\allowbreak{}t\allowbreak{}a\allowbreak{}P\allowbreak{}r\allowbreak{}o\allowbreak{}v\allowbreak{}i\allowbreak{}d\allowbreak{}e\allowbreak{}r} metadata is available for which resource, but it does not subscribe, refresh, resolve nested scopes, or dispatch Commands.

\subsection{Binding Snapshot}\label{record-132-binding-snapshot}

An immutable runtime value created by a provider for one requirement.

A snapshot has:

\begin{itemize}
\tightlist
\item
  provider id
\item
  requirement id
\item
  version
\item
  status
\item
  data when ready
\item
  issues when failed
\item
  freshness facts when stale
\item
  semantic facts for tooling and lower modes
\end{itemize}

\subsection{Command}\label{record-132-command}

A typed intent emitted by a view in response to user input.

Commands say what the user intended. They do not contain hidden provider mutations.

\section{Architecture}\label{record-132-architecture}

\subsection{1. Views Declare Requirements}\label{record-132-views-declare-requirements}

Blocks and views should declare data requirements alongside metadata:

\begin{CodeBlock}
const readerSurfaceBlock = defineBlock({
  metadata: readerSurfaceMetadata,
  data: defineViewData({
    article: requiredResource(articleResource, {
      id: routeParam('articleId'),
    }),
    outline: optionalResource(articleOutlineResource, {
      id: routeParam('articleId'),
    }),
  }),
  commands: {
    selectHeading: commandIntent<{ headingId: string }>('reader.selectHeading'),
  },
  render: ({ data, command }) => readerSurface({
    article: data.require('article'),
    outline: data.get('outline'),
    onSelectHeading: headingId => command('selectHeading', { headingId }),
  }),
});
\end{CodeBlock}

This sketch is not final API. The required properties are:

\begin{itemize}
\tightlist
\item
  data requirements are declared, not discovered from render-time reads
\item
  requirements are named
\item
  requirements are inspectable before rendering
\item
  Commands are declared as intent outputs
\item
  render receives immutable binding data, not provider objects
\end{itemize}

\subsection{2. Active Views Determine Active Bindings}\label{record-132-active-views-determine-active-bindings}

The runtime should resolve data only for views in the active view hierarchy.

When a view becomes active:

\begin{enumerate}
\def\labelenumi{\arabic{enumi}.}
\tightlist
\item
  collect its data requirements
\item
  resolve each requirement against the nearest binding scope
\item
  request or subscribe to provider snapshots
\item
  create a binding frame
\item
  render the view from that immutable frame
\end{enumerate}

When a view becomes inactive:

\begin{enumerate}
\def\labelenumi{\arabic{enumi}.}
\tightlist
\item
  dispose or suspend bindings owned only by that view
\item
  preserve shared provider caches only when the provider owns that policy
\item
  remove binding facts from active runtime introspection
\end{enumerate}

The view tree does not manually subscribe or unsubscribe.

\subsection{3. Providers Normalize Every Backing Source}\label{record-132-providers-normalize-every-backing-source}

Providers hide transport details.

Whether data comes from an event, a message channel, a query, a file, or memory, the provider contract should normalize it into immutable snapshots:

\begin{CodeBlock}
interface BindingSnapshot<Data> {
  readonly providerId: string;
  readonly requirementId: string;
  readonly version: number;
  readonly status: 'ready' | 'loading' | 'empty' | 'stale' | 'error';
  readonly data?: DeepReadonly<Data>;
  readonly issues: readonly BindingIssue[];
  readonly facts: readonly BindingFact[];
}
\end{CodeBlock}

Provider updates do not mutate mounted views. They produce a new snapshot version and trigger a runtime binding invalidation. The next render receives a new immutable binding frame.

DX-034A snapshots accept inert plain snapshot data: null, strings, numbers, booleans, arrays, and enumerable string-keyed plain objects. They reject mutable built-ins and executable values such as \texttt{M\allowbreak{}a\allowbreak{}p}, \texttt{S\allowbreak{}e\allowbreak{}t}, \texttt{D\allowbreak{}a\allowbreak{}t\allowbreak{}e}, typed arrays, functions, accessors, symbol-keyed properties, symbols, and bigint. Snapshot construction copies accepted data before freezing it, so provider-owned or business-owned source objects are not frozen in place.

\subsection{4. Binding Frames Are Immutable}\label{record-132-binding-frames-are-immutable}

Render receives a binding frame, not mutable provider handles.

The frame should support explicit reads:

\begin{CodeBlock}
data.require('article');
data.get('outline');
data.status('comments');
data.issues('reviews');
\end{CodeBlock}

The frame should not support writes.

Any attempt to mutate data from inside a view should be structurally impossible in the public API and rejected at runtime where objects cross a boundary.

\subsection{DX-034A Binding Snapshot And Frame Primitives}\label{record-132-dx-034a-binding-snapshot-and-frame-primitives}

DX-034A lands the first runtime-truth slice in \texttt{@\allowbreak{}f\allowbreak{}l\allowbreak{}y\allowbreak{}i\allowbreak{}n\allowbreak{}g\allowbreak{}r\allowbreak{}o\allowbreak{}b\allowbreak{}o\allowbreak{}t\allowbreak{}s\allowbreak{}/\allowbreak{}b\allowbreak{}i\allowbreak{}j\allowbreak{}o\allowbreak{}u}.

The public core now includes:

\begin{itemize}
\tightlist
\item
  \texttt{D\allowbreak{}a\allowbreak{}t\allowbreak{}a\allowbreak{}R\allowbreak{}e\allowbreak{}q\allowbreak{}u\allowbreak{}i\allowbreak{}r\allowbreak{}e\allowbreak{}m\allowbreak{}e\allowbreak{}n\allowbreak{}t}
\item
  \texttt{P\allowbreak{}r\allowbreak{}o\allowbreak{}v\allowbreak{}i\allowbreak{}d\allowbreak{}e\allowbreak{}r\allowbreak{}I\allowbreak{}d}
\item
  \texttt{R\allowbreak{}e\allowbreak{}q\allowbreak{}u\allowbreak{}i\allowbreak{}r\allowbreak{}e\allowbreak{}m\allowbreak{}e\allowbreak{}n\allowbreak{}t\allowbreak{}I\allowbreak{}d}
\item
  \texttt{B\allowbreak{}i\allowbreak{}n\allowbreak{}d\allowbreak{}i\allowbreak{}n\allowbreak{}g\allowbreak{}S\allowbreak{}t\allowbreak{}a\allowbreak{}t\allowbreak{}u\allowbreak{}s}
\item
  \texttt{B\allowbreak{}i\allowbreak{}n\allowbreak{}d\allowbreak{}i\allowbreak{}n\allowbreak{}g\allowbreak{}S\allowbreak{}n\allowbreak{}a\allowbreak{}p\allowbreak{}s\allowbreak{}h\allowbreak{}o\allowbreak{}t}
\item
  \texttt{B\allowbreak{}i\allowbreak{}n\allowbreak{}d\allowbreak{}i\allowbreak{}n\allowbreak{}g\allowbreak{}F\allowbreak{}r\allowbreak{}a\allowbreak{}m\allowbreak{}e}
\item
  \texttt{B\allowbreak{}i\allowbreak{}n\allowbreak{}d\allowbreak{}i\allowbreak{}n\allowbreak{}g\allowbreak{}I\allowbreak{}s\allowbreak{}s\allowbreak{}u\allowbreak{}e}
\item
  \texttt{B\allowbreak{}i\allowbreak{}n\allowbreak{}d\allowbreak{}i\allowbreak{}n\allowbreak{}g\allowbreak{}F\allowbreak{}a\allowbreak{}c\allowbreak{}t}
\item
  \texttt{C\allowbreak{}o\allowbreak{}m\allowbreak{}m\allowbreak{}a\allowbreak{}n\allowbreak{}d\allowbreak{}I\allowbreak{}n\allowbreak{}t\allowbreak{}e\allowbreak{}n\allowbreak{}t}
\item
  \texttt{d\allowbreak{}e\allowbreak{}f\allowbreak{}i\allowbreak{}n\allowbreak{}e\allowbreak{}D\allowbreak{}a\allowbreak{}t\allowbreak{}a\allowbreak{}R\allowbreak{}e\allowbreak{}q\allowbreak{}u\allowbreak{}i\allowbreak{}r\allowbreak{}e\allowbreak{}m\allowbreak{}e\allowbreak{}n\allowbreak{}t\allowbreak{}(\allowbreak{})}
\item
  \texttt{b\allowbreak{}i\allowbreak{}n\allowbreak{}d\allowbreak{}i\allowbreak{}n\allowbreak{}g\allowbreak{}S\allowbreak{}n\allowbreak{}a\allowbreak{}p\allowbreak{}s\allowbreak{}h\allowbreak{}o\allowbreak{}t\allowbreak{}(\allowbreak{})}
\item
  \texttt{b\allowbreak{}i\allowbreak{}n\allowbreak{}d\allowbreak{}i\allowbreak{}n\allowbreak{}g\allowbreak{}F\allowbreak{}r\allowbreak{}a\allowbreak{}m\allowbreak{}e\allowbreak{}(\allowbreak{})}
\item
  \texttt{c\allowbreak{}o\allowbreak{}m\allowbreak{}m\allowbreak{}a\allowbreak{}n\allowbreak{}d\allowbreak{}I\allowbreak{}n\allowbreak{}t\allowbreak{}e\allowbreak{}n\allowbreak{}t\allowbreak{}(\allowbreak{})}
\end{itemize}

This slice proves the objects, not the full runtime lifecycle:

\begin{CodeBlock}
const article = defineDataRequirement({
  id: 'article',
  resource: 'docs.article',
});

const snapshot = bindingSnapshot({
  providerId: 'docs.articleProvider',
  requirementId: article.id,
  version: 1,
  status: 'ready',
  data: { title: 'DX-034' },
});

const frame = bindingFrame([snapshot]);

frame.require('article');
frame.get('article');
frame.status('article');
frame.issues('article');
\end{CodeBlock}

The \texttt{B\allowbreak{}i\allowbreak{}n\allowbreak{}d\allowbreak{}i\allowbreak{}n\allowbreak{}g\allowbreak{}F\allowbreak{}r\allowbreak{}a\allowbreak{}m\allowbreak{}e} deliberately does not include provider scopes, subscriptions, active hierarchy traversal, schema adapters, command dispatch, or AppShell rendering. Those remain separate runtime layers.

\subsection{DX-034B Provider Scope Contracts}\label{record-132-dx-034b-provider-scope-contracts}

DX-034B lands the explicit provider/scope contract layer in \texttt{@\allowbreak{}f\allowbreak{}l\allowbreak{}y\allowbreak{}i\allowbreak{}n\allowbreak{}g\allowbreak{}r\allowbreak{}o\allowbreak{}b\allowbreak{}o\allowbreak{}t\allowbreak{}s\allowbreak{}/\allowbreak{}b\allowbreak{}i\allowbreak{}j\allowbreak{}o\allowbreak{}u}.

The public core now includes:

\begin{itemize}
\tightlist
\item
  \texttt{D\allowbreak{}a\allowbreak{}t\allowbreak{}a\allowbreak{}P\allowbreak{}r\allowbreak{}o\allowbreak{}v\allowbreak{}i\allowbreak{}d\allowbreak{}e\allowbreak{}r}
\item
  \texttt{D\allowbreak{}a\allowbreak{}t\allowbreak{}a\allowbreak{}P\allowbreak{}r\allowbreak{}o\allowbreak{}v\allowbreak{}i\allowbreak{}d\allowbreak{}e\allowbreak{}r\allowbreak{}I\allowbreak{}n\allowbreak{}p\allowbreak{}u\allowbreak{}t}
\item
  \texttt{P\allowbreak{}r\allowbreak{}o\allowbreak{}v\allowbreak{}i\allowbreak{}d\allowbreak{}e\allowbreak{}r\allowbreak{}S\allowbreak{}c\allowbreak{}o\allowbreak{}p\allowbreak{}e}
\item
  \texttt{P\allowbreak{}r\allowbreak{}o\allowbreak{}v\allowbreak{}i\allowbreak{}d\allowbreak{}e\allowbreak{}r\allowbreak{}S\allowbreak{}c\allowbreak{}o\allowbreak{}p\allowbreak{}e\allowbreak{}E\allowbreak{}n\allowbreak{}t\allowbreak{}r\allowbreak{}y}
\item
  \texttt{P\allowbreak{}r\allowbreak{}o\allowbreak{}v\allowbreak{}i\allowbreak{}d\allowbreak{}e\allowbreak{}r\allowbreak{}S\allowbreak{}c\allowbreak{}o\allowbreak{}p\allowbreak{}e\allowbreak{}O\allowbreak{}p\allowbreak{}t\allowbreak{}i\allowbreak{}o\allowbreak{}n\allowbreak{}s}
\item
  \texttt{d\allowbreak{}e\allowbreak{}f\allowbreak{}i\allowbreak{}n\allowbreak{}e\allowbreak{}D\allowbreak{}a\allowbreak{}t\allowbreak{}a\allowbreak{}P\allowbreak{}r\allowbreak{}o\allowbreak{}v\allowbreak{}i\allowbreak{}d\allowbreak{}e\allowbreak{}r\allowbreak{}(\allowbreak{})}
\item
  \texttt{p\allowbreak{}r\allowbreak{}o\allowbreak{}v\allowbreak{}i\allowbreak{}d\allowbreak{}e\allowbreak{}(\allowbreak{})}
\item
  \texttt{p\allowbreak{}r\allowbreak{}o\allowbreak{}v\allowbreak{}i\allowbreak{}d\allowbreak{}e\allowbreak{}r\allowbreak{}S\allowbreak{}c\allowbreak{}o\allowbreak{}p\allowbreak{}e\allowbreak{}(\allowbreak{})}
\item
  \texttt{i\allowbreak{}s\allowbreak{}D\allowbreak{}a\allowbreak{}t\allowbreak{}a\allowbreak{}P\allowbreak{}r\allowbreak{}o\allowbreak{}v\allowbreak{}i\allowbreak{}d\allowbreak{}e\allowbreak{}r\allowbreak{}(\allowbreak{})}
\item
  \texttt{P\allowbreak{}r\allowbreak{}o\allowbreak{}v\allowbreak{}i\allowbreak{}d\allowbreak{}e\allowbreak{}r\allowbreak{}R\allowbreak{}e\allowbreak{}s\allowbreak{}o\allowbreak{}l\allowbreak{}u\allowbreak{}t\allowbreak{}i\allowbreak{}o\allowbreak{}n}
\item
  \texttt{P\allowbreak{}r\allowbreak{}o\allowbreak{}v\allowbreak{}i\allowbreak{}d\allowbreak{}e\allowbreak{}r\allowbreak{}R\allowbreak{}e\allowbreak{}s\allowbreak{}o\allowbreak{}l\allowbreak{}u\allowbreak{}t\allowbreak{}i\allowbreak{}o\allowbreak{}n\allowbreak{}S\allowbreak{}t\allowbreak{}a\allowbreak{}t\allowbreak{}u\allowbreak{}s}
\item
  \texttt{r\allowbreak{}e\allowbreak{}s\allowbreak{}o\allowbreak{}l\allowbreak{}v\allowbreak{}e\allowbreak{}P\allowbreak{}r\allowbreak{}o\allowbreak{}v\allowbreak{}i\allowbreak{}d\allowbreak{}e\allowbreak{}r\allowbreak{}R\allowbreak{}e\allowbreak{}q\allowbreak{}u\allowbreak{}i\allowbreak{}r\allowbreak{}e\allowbreak{}m\allowbreak{}e\allowbreak{}n\allowbreak{}t\allowbreak{}(\allowbreak{})}
\item
  \texttt{r\allowbreak{}e\allowbreak{}s\allowbreak{}o\allowbreak{}l\allowbreak{}v\allowbreak{}e\allowbreak{}P\allowbreak{}r\allowbreak{}o\allowbreak{}v\allowbreak{}i\allowbreak{}d\allowbreak{}e\allowbreak{}r\allowbreak{}R\allowbreak{}e\allowbreak{}q\allowbreak{}u\allowbreak{}i\allowbreak{}r\allowbreak{}e\allowbreak{}m\allowbreak{}e\allowbreak{}n\allowbreak{}t\allowbreak{}s\allowbreak{}(\allowbreak{})}
\item
  \texttt{B\allowbreak{}i\allowbreak{}n\allowbreak{}d\allowbreak{}i\allowbreak{}n\allowbreak{}g\allowbreak{}F\allowbreak{}r\allowbreak{}a\allowbreak{}m\allowbreak{}e\allowbreak{}A\allowbreak{}s\allowbreak{}s\allowbreak{}e\allowbreak{}m\allowbreak{}b\allowbreak{}l\allowbreak{}y}
\item
  \texttt{b\allowbreak{}i\allowbreak{}n\allowbreak{}d\allowbreak{}i\allowbreak{}n\allowbreak{}g\allowbreak{}F\allowbreak{}r\allowbreak{}a\allowbreak{}m\allowbreak{}e\allowbreak{}F\allowbreak{}r\allowbreak{}o\allowbreak{}m\allowbreak{}S\allowbreak{}n\allowbreak{}a\allowbreak{}p\allowbreak{}s\allowbreak{}h\allowbreak{}o\allowbreak{}t\allowbreak{}s\allowbreak{}(\allowbreak{})}
\item
  \texttt{V\allowbreak{}i\allowbreak{}e\allowbreak{}w\allowbreak{}D\allowbreak{}a\allowbreak{}t\allowbreak{}a\allowbreak{}C\allowbreak{}o\allowbreak{}n\allowbreak{}t\allowbreak{}r\allowbreak{}a\allowbreak{}c\allowbreak{}t}
\item
  \texttt{V\allowbreak{}i\allowbreak{}e\allowbreak{}w\allowbreak{}D\allowbreak{}a\allowbreak{}t\allowbreak{}a\allowbreak{}R\allowbreak{}e\allowbreak{}q\allowbreak{}u\allowbreak{}i\allowbreak{}r\allowbreak{}e\allowbreak{}m\allowbreak{}e\allowbreak{}n\allowbreak{}t\allowbreak{}E\allowbreak{}n\allowbreak{}t\allowbreak{}r\allowbreak{}y}
\item
  \texttt{d\allowbreak{}e\allowbreak{}f\allowbreak{}i\allowbreak{}n\allowbreak{}e\allowbreak{}V\allowbreak{}i\allowbreak{}e\allowbreak{}w\allowbreak{}D\allowbreak{}a\allowbreak{}t\allowbreak{}a\allowbreak{}(\allowbreak{})}
\end{itemize}

This slice proves local provider availability and one-scope resolution, not hierarchical provider resolution or subscription lifecycle:

\begin{CodeBlock}
const article = defineDataRequirement({
  id: 'article',
  resource: 'docs.article',
});

const readerData = defineViewData({
  id: 'reader.data',
  requirements: [{ name: 'article', requirement: article }],
});

const articleProvider = defineDataProvider({
  id: 'docs.articleProvider',
  resource: article.resource,
});

const providers = providerScope([provide(articleProvider)], {
  id: 'docs.appShell',
});

providers.has(article.resource);
providers.get(article.resource);
providers.resources();

const resolution = resolveProviderRequirement(article, providers);
resolution.status;

const assembled = bindingFrameFromSnapshots({
  resolutions: [resolution],
  snapshots: [
    bindingSnapshot({
      providerId: 'docs.articleProvider',
      requirementId: 'article',
      version: 1,
      status: 'ready',
      data: { title: 'DX-034' },
    }),
  ],
});
assembled.frame.require('article');
\end{CodeBlock}

\texttt{P\allowbreak{}r\allowbreak{}o\allowbreak{}v\allowbreak{}i\allowbreak{}d\allowbreak{}e\allowbreak{}r\allowbreak{}S\allowbreak{}c\allowbreak{}o\allowbreak{}p\allowbreak{}e} rejects duplicate resources and duplicate provider ids inside a single scope. It also rejects loose provider-shaped objects, so importing a module cannot silently register a provider or bypass the constructor contract. \texttt{V\allowbreak{}i\allowbreak{}e\allowbreak{}w\allowbreak{}D\allowbreak{}a\allowbreak{}t\allowbreak{}a\allowbreak{}C\allowbreak{}o\allowbreak{}n\allowbreak{}t\allowbreak{}r\allowbreak{}a\allowbreak{}c\allowbreak{}t} names the data a view needs before any provider binding occurs. Resolution results are frozen metadata. Missing required providers become deterministic issues; missing optional providers are inspectable but not errors. Frame assembly reports missing or provider-mismatched snapshots as deterministic issues while keeping provider reads outside the render frame.

\subsection{DX-034C AppShell Composition Contract}\label{record-132-dx-034c-appshell-composition-contract}

DX-034C lands the first structural AppShell composition contract in \texttt{@\allowbreak{}f\allowbreak{}l\allowbreak{}y\allowbreak{}i\allowbreak{}n\allowbreak{}g\allowbreak{}r\allowbreak{}o\allowbreak{}b\allowbreak{}o\allowbreak{}t\allowbreak{}s\allowbreak{}/\allowbreak{}b\allowbreak{}i\allowbreak{}j\allowbreak{}o\allowbreak{}u}.

The public core now includes:

\begin{itemize}
\tightlist
\item
  \texttt{A\allowbreak{}p\allowbreak{}p\allowbreak{}S\allowbreak{}h\allowbreak{}e\allowbreak{}l\allowbreak{}l\allowbreak{}C\allowbreak{}o\allowbreak{}m\allowbreak{}p\allowbreak{}o\allowbreak{}s\allowbreak{}i\allowbreak{}t\allowbreak{}i\allowbreak{}o\allowbreak{}n}
\item
  \texttt{A\allowbreak{}p\allowbreak{}p\allowbreak{}S\allowbreak{}h\allowbreak{}e\allowbreak{}l\allowbreak{}l\allowbreak{}C\allowbreak{}o\allowbreak{}m\allowbreak{}p\allowbreak{}o\allowbreak{}s\allowbreak{}i\allowbreak{}t\allowbreak{}i\allowbreak{}o\allowbreak{}n\allowbreak{}I\allowbreak{}n\allowbreak{}p\allowbreak{}u\allowbreak{}t}
\item
  \texttt{A\allowbreak{}p\allowbreak{}p\allowbreak{}S\allowbreak{}h\allowbreak{}e\allowbreak{}l\allowbreak{}l\allowbreak{}S\allowbreak{}l\allowbreak{}o\allowbreak{}t}
\item
  \texttt{A\allowbreak{}p\allowbreak{}p\allowbreak{}S\allowbreak{}h\allowbreak{}e\allowbreak{}l\allowbreak{}l\allowbreak{}S\allowbreak{}l\allowbreak{}o\allowbreak{}t\allowbreak{}C\allowbreak{}o\allowbreak{}n\allowbreak{}t\allowbreak{}e\allowbreak{}n\allowbreak{}t}
\item
  \texttt{A\allowbreak{}p\allowbreak{}p\allowbreak{}S\allowbreak{}h\allowbreak{}e\allowbreak{}l\allowbreak{}l\allowbreak{}S\allowbreak{}l\allowbreak{}o\allowbreak{}t\allowbreak{}I\allowbreak{}d}
\item
  \texttt{A\allowbreak{}p\allowbreak{}p\allowbreak{}S\allowbreak{}h\allowbreak{}e\allowbreak{}l\allowbreak{}l\allowbreak{}S\allowbreak{}l\allowbreak{}o\allowbreak{}t\allowbreak{}s}
\item
  \texttt{d\allowbreak{}e\allowbreak{}f\allowbreak{}i\allowbreak{}n\allowbreak{}e\allowbreak{}A\allowbreak{}p\allowbreak{}p\allowbreak{}S\allowbreak{}h\allowbreak{}e\allowbreak{}l\allowbreak{}l\allowbreak{}C\allowbreak{}o\allowbreak{}m\allowbreak{}p\allowbreak{}o\allowbreak{}s\allowbreak{}i\allowbreak{}t\allowbreak{}i\allowbreak{}o\allowbreak{}n\allowbreak{}(\allowbreak{})}
\item
  \texttt{i\allowbreak{}s\allowbreak{}A\allowbreak{}p\allowbreak{}p\allowbreak{}S\allowbreak{}h\allowbreak{}e\allowbreak{}l\allowbreak{}l\allowbreak{}C\allowbreak{}o\allowbreak{}m\allowbreak{}p\allowbreak{}o\allowbreak{}s\allowbreak{}i\allowbreak{}t\allowbreak{}i\allowbreak{}o\allowbreak{}n\allowbreak{}(\allowbreak{})}
\end{itemize}

This slice proves semantic shell slots and inspectability. It still does not render AppShell, resolve hierarchical provider scopes, subscribe, refresh, walk the active hierarchy, or dispatch commands:

\begin{CodeBlock}
const docsShell = defineAppShellComposition({
  id: 'docs.shell',
  providers,
  slots: {
    navigation: navigationBlock,
    content: readerSurfaceBlock,
    inspector: inspectorPanelBlock,
  },
});

docsShell.slotIds();
docsShell.slot('content');
docsShell.dataContracts();
docsShell.commandIntents();
docsShell.providerScope();
\end{CodeBlock}

Slots are logical regions: \texttt{n\allowbreak{}a\allowbreak{}v\allowbreak{}i\allowbreak{}g\allowbreak{}a\allowbreak{}t\allowbreak{}i\allowbreak{}o\allowbreak{}n}, \texttt{c\allowbreak{}o\allowbreak{}n\allowbreak{}t\allowbreak{}e\allowbreak{}n\allowbreak{}t}, \texttt{i\allowbreak{}n\allowbreak{}s\allowbreak{}p\allowbreak{}e\allowbreak{}c\allowbreak{}t\allowbreak{}o\allowbreak{}r}, \texttt{s\allowbreak{}t\allowbreak{}a\allowbreak{}t\allowbreak{}u\allowbreak{}s}, and \texttt{o\allowbreak{}v\allowbreak{}e\allowbreak{}r\allowbreak{}l\allowbreak{}a\allowbreak{}y\allowbreak{}s}. Physical slots such as \texttt{l\allowbreak{}e\allowbreak{}f\allowbreak{}t\allowbreak{}N\allowbreak{}a\allowbreak{}v} are rejected at construction time. Slot content must be a block returned by \texttt{d\allowbreak{}e\allowbreak{}f\allowbreak{}i\allowbreak{}n\allowbreak{}e\allowbreak{}B\allowbreak{}l\allowbreak{}o\allowbreak{}c\allowbreak{}k\allowbreak{}(\allowbreak{})} or nested arrays of such blocks. The composition exposes the nested blocks' declared view data contracts and command intents for tooling without adding provider handles or render-time mutation paths.

\subsection{DX-034D Active Binding Lifecycle}\label{record-132-dx-034d-active-binding-lifecycle}

DX-034D formalizes the active binding lifecycle before rendered AppShell or runtime provider subscriptions begin.

Core invariant:

\begin{CodeBlock}
Active hierarchy owns which bindings are active.
Providers publish snapshots.
Runtime assembles new frames.
Views render frames.
Commands leave views as intent.
Business logic owns change.
\end{CodeBlock}

This slice is design-only. It defines lifecycle ownership and vocabulary; it does not implement provider subscriptions, active hierarchy traversal, command dispatch, rendered AppShell, schema binding, cache retention, or DOGFOOD integration.

\subsubsection{Lifecycle Ownership}\label{record-132-lifecycle-ownership}

The active view hierarchy owns binding lifetime.

A data requirement becomes active when all of these are true:

\begin{enumerate}
\def\labelenumi{\arabic{enumi}.}
\tightlist
\item
  a block or view declares the requirement through an inspectable data contract
\item
  that block or view is present in the active hierarchy
\item
  the runtime resolves the requirement against an explicit provider scope
\item
  the runtime creates binding ownership for that active view node
\end{enumerate}

Views never choose their own binding state. Views do not subscribe, refresh, dispose, retain cache, or mutate provider state. The runtime owns lifecycle transitions and provider subscription policy.

\subsubsection{Lifecycle Vocabulary}\label{record-132-lifecycle-vocabulary}

Use the smallest useful lifecycle vocabulary first:

\begin{itemize}
\tightlist
\item
  \texttt{a\allowbreak{}c\allowbreak{}t\allowbreak{}i\allowbreak{}v\allowbreak{}e}: the requirement is owned by the active hierarchy; the runtime may hold a provider subscription; provider snapshots can invalidate the current frame.
\item
  \texttt{s\allowbreak{}u\allowbreak{}s\allowbreak{}p\allowbreak{}e\allowbreak{}n\allowbreak{}d\allowbreak{}e\allowbreak{}d}: the requirement is not currently rendering, but runtime/provider policy may retain resumable state or cache. Views do not decide what is retained.
\item
  \texttt{d\allowbreak{}i\allowbreak{}s\allowbreak{}p\allowbreak{}o\allowbreak{}s\allowbreak{}e\allowbreak{}d}: the requirement ownership is released. The runtime no longer expects provider updates for that view binding.
\end{itemize}

Provider cache retention is provider/runtime policy, not view policy. DX-034D names lifecycle states but does not implement cache retention.

\subsubsection{Provider Subscription Ownership}\label{record-132-provider-subscription-ownership}

Provider subscriptions are runtime lifecycle. Binding frames are immutable render inputs.

Providers may subscribe to channels, messages, queries, files, databases, MCP, or in-memory business state internally, but render receives only immutable snapshots assembled into a \texttt{B\allowbreak{}i\allowbreak{}n\allowbreak{}d\allowbreak{}i\allowbreak{}n\allowbreak{}g\allowbreak{}F\allowbreak{}r\allowbreak{}a\allowbreak{}m\allowbreak{}e}. A view never receives provider handles, subscription handles, refresh methods, mutable stores, or dispatcher callbacks.

\subsubsection{Provider Update Flow}\label{record-132-provider-update-flow}

Provider updates do not mutate rendered views or existing frames.

The lifecycle flow is:

% Presentation of original arrow notation. See flow-manifest.json.
\begin{center}
\begin{tikzpicture}[book flow node/.style={draw=BookHeading,line width=.65pt,fill=BookPanel,rounded corners=1.5pt,align=center,inner sep=5pt,minimum height=10mm,font=\fontsize{9}{11}\selectfont,execute at begin node={\hyphenpenalty=10000\exhyphenpenalty=10000}},book flow edge/.style={-{Latex[length=2mm]},draw=BookHeading,line width=.7pt}]
\node[book flow node,text width=96mm] (n0) at (0,0.0) {provider publishes snapshot};
\node[book flow node,text width=96mm] (n1) at (0,-1.55) {runtime records binding invalidation};
\node[book flow node,text width=96mm] (n2) at (0,-3.1) {runtime assembles next Binding\allowbreak{}Frame};
\node[book flow node,text width=96mm] (n3) at (0,-4.65) {active view renders the next frame};
\draw[book flow edge] (n0) -- (n1);
\draw[book flow edge] (n1) -- (n2);
\draw[book flow edge] (n2) -- (n3);
\end{tikzpicture}
\end{center}


Invalidation belongs to runtime binding ownership, not to rendered cells. Static, pipe, and accessible modes must be able to observe lifecycle facts such as active, suspended, disposed, loading, stale, empty, error, and invalidated without parsing terminal output.

\subsubsection{Command Flow}\label{record-132-command-flow}

Commands leave views as intent records.

Views may emit declared command intents, but they do not dispatch business logic, call provider methods, or mutate provider inputs. Runtime command integration receives command intent records and hands them to business logic. Business logic handles Commands, updates state or provider inputs, and providers eventually publish new snapshots.

\subsubsection{Active Lifecycle Questions}\label{record-132-active-lifecycle-questions}

DX-034D answers:

\begin{itemize}
\tightlist
\item
  When does a requirement become active?
\item
  Who owns the binding while a view is active?
\item
  What happens when a view leaves the active hierarchy?
\item
  What is suspended versus disposed?
\item
  Who owns provider subscriptions?
\item
  How does a provider update become a new \texttt{B\allowbreak{}i\allowbreak{}n\allowbreak{}d\allowbreak{}i\allowbreak{}n\allowbreak{}g\allowbreak{}F\allowbreak{}r\allowbreak{}a\allowbreak{}m\allowbreak{}e}?
\item
  How do Commands leave views?
\item
  Who dispatches Commands to business logic?
\item
  What lifecycle facts are inspectable?
\item
  What lower-mode facts survive static, pipe, and accessible output?
\end{itemize}

\subsection{DX-034E Binding Lifecycle Primitives}\label{record-132-dx-034e-binding-lifecycle-primitives}

DX-034E implements lifecycle as transition algebra, not as a runtime manager.

The primitives answer what state a binding is in, who owns it, what transition produced the next record, what invalidation happened, and what lifecycle facts tooling can inspect.

They do not subscribe, refresh, dispatch, traverse the active view tree, render AppShell, retain caches, or bind schemas.

The public vocabulary stays deliberately small:

\begin{itemize}
\tightlist
\item
  \texttt{B\allowbreak{}i\allowbreak{}n\allowbreak{}d\allowbreak{}i\allowbreak{}n\allowbreak{}g\allowbreak{}L\allowbreak{}i\allowbreak{}f\allowbreak{}e\allowbreak{}c\allowbreak{}y\allowbreak{}c\allowbreak{}l\allowbreak{}e\allowbreak{}O\allowbreak{}w\allowbreak{}n\allowbreak{}e\allowbreak{}r}: the active-view, block, AppShell, or runtime owner responsible for one binding lifetime.
\item
  \texttt{B\allowbreak{}i\allowbreak{}n\allowbreak{}d\allowbreak{}i\allowbreak{}n\allowbreak{}g\allowbreak{}L\allowbreak{}i\allowbreak{}f\allowbreak{}e\allowbreak{}c\allowbreak{}y\allowbreak{}c\allowbreak{}l\allowbreak{}e\allowbreak{}R\allowbreak{}e\allowbreak{}c\allowbreak{}o\allowbreak{}r\allowbreak{}d}: an immutable record for one owner and one requirement.
\item
  \texttt{B\allowbreak{}i\allowbreak{}n\allowbreak{}d\allowbreak{}i\allowbreak{}n\allowbreak{}g\allowbreak{}L\allowbreak{}i\allowbreak{}f\allowbreak{}e\allowbreak{}c\allowbreak{}y\allowbreak{}c\allowbreak{}l\allowbreak{}e\allowbreak{}S\allowbreak{}t\allowbreak{}a\allowbreak{}t\allowbreak{}e}: \texttt{a\allowbreak{}c\allowbreak{}t\allowbreak{}i\allowbreak{}v\allowbreak{}e}, \texttt{s\allowbreak{}u\allowbreak{}s\allowbreak{}p\allowbreak{}e\allowbreak{}n\allowbreak{}d\allowbreak{}e\allowbreak{}d}, and \texttt{d\allowbreak{}i\allowbreak{}s\allowbreak{}p\allowbreak{}o\allowbreak{}s\allowbreak{}e\allowbreak{}d}.
\item
  \texttt{B\allowbreak{}i\allowbreak{}n\allowbreak{}d\allowbreak{}i\allowbreak{}n\allowbreak{}g\allowbreak{}L\allowbreak{}i\allowbreak{}f\allowbreak{}e\allowbreak{}c\allowbreak{}y\allowbreak{}c\allowbreak{}l\allowbreak{}e\allowbreak{}T\allowbreak{}r\allowbreak{}a\allowbreak{}n\allowbreak{}s\allowbreak{}i\allowbreak{}t\allowbreak{}i\allowbreak{}o\allowbreak{}n}: the immutable record of an activation, suspension, disposal, or invalidation transition.
\item
  \texttt{B\allowbreak{}i\allowbreak{}n\allowbreak{}d\allowbreak{}i\allowbreak{}n\allowbreak{}g\allowbreak{}I\allowbreak{}n\allowbreak{}v\allowbreak{}a\allowbreak{}l\allowbreak{}i\allowbreak{}d\allowbreak{}a\allowbreak{}t\allowbreak{}i\allowbreak{}o\allowbreak{}n}: an immutable fact that a provider update, scope change, manual runtime action, owner suspension, or owner disposal means a later runtime layer should assemble a new frame.
\end{itemize}

Allowed transitions are:

% Presentation of original arrow notation. See flow-manifest.json.
\begin{center}
\begin{tikzpicture}[book flow node/.style={draw=BookHeading,line width=.65pt,fill=BookPanel,rounded corners=1.5pt,align=center,inner sep=5pt,minimum height=10mm,font=\fontsize{9}{11}\selectfont,execute at begin node={\hyphenpenalty=10000\exhyphenpenalty=10000}},book flow edge/.style={-{Latex[length=2mm]},draw=BookHeading,line width=.7pt}]
\node[book flow node,text width=36mm] (n0) at (0,0) {active};
\node[book flow node,text width=36mm] (n1) at (6,0) {suspended};
\node[book flow node,text width=36mm] (n2) at (3,-4) {disposed};
\draw[book flow edge] (n0) to[bend left=18]  (n1);
\draw[book flow edge] (n0) -- (n2);
\draw[book flow edge] (n1) to[bend left=18]  (n0);
\draw[book flow edge] (n1) -- (n2);
\draw[book flow edge] (n0) to[loop above]  node[fill=white,inner sep=2pt,align=center,font=\scriptsize] {by invalidation only} (n0);
\draw[book flow edge] (n1) to[loop above]  node[fill=white,inner sep=2pt,align=center,font=\scriptsize] {by invalidation only} (n1);
\end{tikzpicture}
\end{center}


Disposed bindings cannot be activated, suspended, or invalidated. A disposed binding means ownership was released; later activation requires a new lifecycle record and a new ownership event.

Provider updates create invalidation records. They do not mutate existing frames, fetch data, call providers, or dispatch Commands:

% Presentation of original arrow notation. See flow-manifest.json.
\begin{center}
\begin{tikzpicture}[book flow node/.style={draw=BookHeading,line width=.65pt,fill=BookPanel,rounded corners=1.5pt,align=center,inner sep=5pt,minimum height=10mm,font=\fontsize{9}{11}\selectfont,execute at begin node={\hyphenpenalty=10000\exhyphenpenalty=10000}},book flow edge/.style={-{Latex[length=2mm]},draw=BookHeading,line width=.7pt}]
\node[book flow node,text width=96mm] (n0) at (0,0.0) {provider publishes snapshot};
\node[book flow node,text width=96mm] (n1) at (0,-1.55) {runtime creates Binding\allowbreak{}Invalidation};
\node[book flow node,text width=96mm] (n2) at (0,-3.1) {runtime later assembles next Binding\allowbreak{}Frame};
\draw[book flow edge] (n0) -- (n1);
\draw[book flow edge] (n1) -- (n2);
\end{tikzpicture}
\end{center}


Lifecycle records are immutable facts, not managers. They expose no provider handles, subscription handles, refresh methods, command dispatch callbacks, mutable stores, or render hooks.

\subsection{DX-034F Active Binding Collection}\label{record-132-dx-034f-active-binding-collection}

DX-034F collects declared active bindings before provider runtime integration.

The collection answers:

\begin{CodeBlock}
given declared data requirements and explicit lifecycle owners,
which owner/requirement pairs are active?
\end{CodeBlock}

It uses \texttt{B\allowbreak{}i\allowbreak{}n\allowbreak{}d\allowbreak{}i\allowbreak{}n\allowbreak{}g\allowbreak{}L\allowbreak{}i\allowbreak{}f\allowbreak{}e\allowbreak{}c\allowbreak{}y\allowbreak{}c\allowbreak{}l\allowbreak{}e\allowbreak{}O\allowbreak{}w\allowbreak{}n\allowbreak{}e\allowbreak{}r} from DX-034E. It does not introduce another ownership vocabulary. An active binding entry binds:

\begin{itemize}
\tightlist
\item
  one explicit \texttt{B\allowbreak{}i\allowbreak{}n\allowbreak{}d\allowbreak{}i\allowbreak{}n\allowbreak{}g\allowbreak{}L\allowbreak{}i\allowbreak{}f\allowbreak{}e\allowbreak{}c\allowbreak{}y\allowbreak{}c\allowbreak{}l\allowbreak{}e\allowbreak{}O\allowbreak{}w\allowbreak{}n\allowbreak{}e\allowbreak{}r}
\item
  one declared \texttt{D\allowbreak{}a\allowbreak{}t\allowbreak{}a\allowbreak{}R\allowbreak{}e\allowbreak{}q\allowbreak{}u\allowbreak{}i\allowbreak{}r\allowbreak{}e\allowbreak{}m\allowbreak{}e\allowbreak{}n\allowbreak{}t}
\item
  optional provider id metadata
\end{itemize}

Collections can produce active \texttt{B\allowbreak{}i\allowbreak{}n\allowbreak{}d\allowbreak{}i\allowbreak{}n\allowbreak{}g\allowbreak{}L\allowbreak{}i\allowbreak{}f\allowbreak{}e\allowbreak{}c\allowbreak{}y\allowbreak{}c\allowbreak{}l\allowbreak{}e\allowbreak{}R\allowbreak{}e\allowbreak{}c\allowbreak{}o\allowbreak{}r\allowbreak{}d} values from those entries so the later runtime layer has stable ownership facts to consume:

\begin{CodeBlock}
const collection = activeBindingCollection([
  activeBindingEntry({
    owner: defineBindingLifecycleOwner({
      id: 'docs.shell',
      kind: 'app-shell',
    }),
    requirement: articleRequirement,
    providerId: 'docs.articleProvider',
  }),
]);

collection.entries();
collection.requirements();
collection.owners();
collection.lifecycleRecords();
\end{CodeBlock}

DX-034F also allows declared \texttt{V\allowbreak{}i\allowbreak{}e\allowbreak{}w\allowbreak{}D\allowbreak{}a\allowbreak{}t\allowbreak{}a\allowbreak{}C\allowbreak{}o\allowbreak{}n\allowbreak{}t\allowbreak{}r\allowbreak{}a\allowbreak{}c\allowbreak{}t} values to be collected without rendering blocks or views:

\begin{CodeBlock}
const collection = collectActiveBindings({
  contracts: [{
    owner,
    contract: readerSurfaceBlock.data,
    providerIds: [
      { requirementId: 'article', providerId: 'docs.articleProvider' },
    ],
  }],
});
\end{CodeBlock}

Collection inspects declarations, not rendered output. It does not subscribe, refresh, dispatch, render, cache, resolve provider scopes, traverse the real TUI runtime tree, bind schemas, or integrate DOGFOOD. Provider ids are metadata only; provider handles are not exposed.

\subsection{DX-034G Active View Binding Collection}\label{record-132-dx-034g-active-view-binding-collection}

DX-034G connects the pure active-binding collection contract to the existing runtime view-stack model in \texttt{@\allowbreak{}f\allowbreak{}l\allowbreak{}y\allowbreak{}i\allowbreak{}n\allowbreak{}g\allowbreak{}r\allowbreak{}o\allowbreak{}b\allowbreak{}o\allowbreak{}t\allowbreak{}s\allowbreak{}/\allowbreak{}b\allowbreak{}i\allowbreak{}j\allowbreak{}o\allowbreak{}u\allowbreak{}-\allowbreak{}t\allowbreak{}u\allowbreak{}i}.

The runtime view stack determines which layers contribute active bindings:

\begin{itemize}
\tightlist
\item
  the top layer is active
\item
  layers below remain active only while each layer above them has \texttt{b\allowbreak{}l\allowbreak{}o\allowbreak{}c\allowbreak{}k\allowbreak{}s\allowbreak{}B\allowbreak{}e\allowbreak{}l\allowbreak{}o\allowbreak{}w\allowbreak{}:\allowbreak{} \allowbreak{}f\allowbreak{}a\allowbreak{}l\allowbreak{}s\allowbreak{}e}
\item
  a layer with \texttt{b\allowbreak{}l\allowbreak{}o\allowbreak{}c\allowbreak{}k\allowbreak{}s\allowbreak{}B\allowbreak{}e\allowbreak{}l\allowbreak{}o\allowbreak{}w\allowbreak{}:\allowbreak{} \allowbreak{}t\allowbreak{}r\allowbreak{}u\allowbreak{}e} blocks lower layers from active binding ownership
\end{itemize}

Runtime binding collection still inspects declarations, not rendered output. Layers expose branded \texttt{R\allowbreak{}u\allowbreak{}n\allowbreak{}t\allowbreak{}i\allowbreak{}m\allowbreak{}e\allowbreak{}V\allowbreak{}i\allowbreak{}e\allowbreak{}w\allowbreak{}B\allowbreak{}i\allowbreak{}n\allowbreak{}d\allowbreak{}i\allowbreak{}n\allowbreak{}g\allowbreak{}S\allowbreak{}o\allowbreak{}u\allowbreak{}r\allowbreak{}c\allowbreak{}e} values, and \texttt{c\allowbreak{}o\allowbreak{}l\allowbreak{}l\allowbreak{}e\allowbreak{}c\allowbreak{}t\allowbreak{}R\allowbreak{}u\allowbreak{}n\allowbreak{}t\allowbreak{}i\allowbreak{}m\allowbreak{}e\allowbreak{}V\allowbreak{}i\allowbreak{}e\allowbreak{}w\allowbreak{}B\allowbreak{}i\allowbreak{}n\allowbreak{}d\allowbreak{}i\allowbreak{}n\allowbreak{}g\allowbreak{}s\allowbreak{}(\allowbreak{})} converts active sources into an \texttt{A\allowbreak{}c\allowbreak{}t\allowbreak{}i\allowbreak{}v\allowbreak{}e\allowbreak{}B\allowbreak{}i\allowbreak{}n\allowbreak{}d\allowbreak{}i\allowbreak{}n\allowbreak{}g\allowbreak{}C\allowbreak{}o\allowbreak{}l\allowbreak{}l\allowbreak{}e\allowbreak{}c\allowbreak{}t\allowbreak{}i\allowbreak{}o\allowbreak{}n}.

Runtime binding sources normalize, copy, and freeze provider assignment metadata at construction so later caller mutation cannot change active binding collection results.

DX-034G does not subscribe, refresh, dispatch, render, cache, resolve provider handles, bind schemas, or integrate DOGFOOD. It only proves that the active view hierarchy can identify active owner/requirement pairs.

\subsection{DX-034H Provider Update Invalidation Flow}\label{record-132-dx-034h-provider-update-invalidation-flow}

DX-034H proves the provider-update half of the loop without implementing a provider manager.

\texttt{b\allowbreak{}i\allowbreak{}n\allowbreak{}d\allowbreak{}i\allowbreak{}n\allowbreak{}g\allowbreak{}F\allowbreak{}r\allowbreak{}a\allowbreak{}m\allowbreak{}e\allowbreak{}U\allowbreak{}p\allowbreak{}d\allowbreak{}a\allowbreak{}t\allowbreak{}e\allowbreak{}F\allowbreak{}r\allowbreak{}o\allowbreak{}m\allowbreak{}S\allowbreak{}n\allowbreak{}a\allowbreak{}p\allowbreak{}s\allowbreak{}h\allowbreak{}o\allowbreak{}t\allowbreak{}s\allowbreak{}(\allowbreak{})} takes:

\begin{itemize}
\tightlist
\item
  an \texttt{A\allowbreak{}c\allowbreak{}t\allowbreak{}i\allowbreak{}v\allowbreak{}e\allowbreak{}B\allowbreak{}i\allowbreak{}n\allowbreak{}d\allowbreak{}i\allowbreak{}n\allowbreak{}g\allowbreak{}C\allowbreak{}o\allowbreak{}l\allowbreak{}l\allowbreak{}e\allowbreak{}c\allowbreak{}t\allowbreak{}i\allowbreak{}o\allowbreak{}n}
\item
  an explicit \texttt{P\allowbreak{}r\allowbreak{}o\allowbreak{}v\allowbreak{}i\allowbreak{}d\allowbreak{}e\allowbreak{}r\allowbreak{}S\allowbreak{}c\allowbreak{}o\allowbreak{}p\allowbreak{}e}
\item
  immutable \texttt{B\allowbreak{}i\allowbreak{}n\allowbreak{}d\allowbreak{}i\allowbreak{}n\allowbreak{}g\allowbreak{}S\allowbreak{}n\allowbreak{}a\allowbreak{}p\allowbreak{}s\allowbreak{}h\allowbreak{}o\allowbreak{}t} values
\item
  optional current lifecycle records
\end{itemize}

It resolves the collection's requirements against the scope, assembles the next immutable \texttt{B\allowbreak{}i\allowbreak{}n\allowbreak{}d\allowbreak{}i\allowbreak{}n\allowbreak{}g\allowbreak{}F\allowbreak{}r\allowbreak{}a\allowbreak{}m\allowbreak{}e}, and records provider updates as lifecycle invalidations. It never mutates the previous frame or lifecycle records. Snapshots whose provider id and version already appear in a lifecycle record's provider-update invalidations do not create duplicate invalidations or lifecycle churn. Explicit provider assignments carried by active binding entries are authoritative: a scope-resolved provider that conflicts with an explicit assignment becomes a deterministic \texttt{p\allowbreak{}r\allowbreak{}o\allowbreak{}v\allowbreak{}i\allowbreak{}d\allowbreak{}e\allowbreak{}r\allowbreak{}.\allowbreak{}a\allowbreak{}s\allowbreak{}s\allowbreak{}i\allowbreak{}g\allowbreak{}n\allowbreak{}m\allowbreak{}e\allowbreak{}n\allowbreak{}t\allowbreak{}-\allowbreak{}m\allowbreak{}i\allowbreak{}s\allowbreak{}m\allowbreak{}a\allowbreak{}t\allowbreak{}c\allowbreak{}h} issue rather than data in the frame.

% Presentation of original arrow notation. See flow-manifest.json.
\begin{center}
\begin{tikzpicture}[book flow node/.style={draw=BookHeading,line width=.65pt,fill=BookPanel,rounded corners=1.5pt,align=center,inner sep=5pt,minimum height=10mm,font=\fontsize{9}{11}\selectfont,execute at begin node={\hyphenpenalty=10000\exhyphenpenalty=10000}},book flow edge/.style={-{Latex[length=2mm]},draw=BookHeading,line width=.7pt}]
\node[book flow node,text width=96mm] (n0) at (0,0.0) {provider publishes snapshot};
\node[book flow node,text width=96mm] (n1) at (0,-1.55) {binding lifecycle records receive invalidation facts};
\node[book flow node,text width=96mm] (n2) at (0,-3.1) {runtime assembles next immutable Binding\allowbreak{}Frame};
\node[book flow node,text width=96mm] (n3) at (0,-4.65) {views later render the new frame};
\draw[book flow edge] (n0) -- (n1);
\draw[book flow edge] (n1) -- (n2);
\draw[book flow edge] (n2) -- (n3);
\end{tikzpicture}
\end{center}


DX-034H does not fetch data, subscribe, refresh providers, retain caches, walk the active hierarchy, render AppShell, or dispatch Commands.

\subsection{DX-034I Command Intent Dispatch Proof}\label{record-132-dx-034i-command-intent-dispatch-proof}

DX-034I proves the command half of the loop without putting business logic in views.

Views emit branded \texttt{R\allowbreak{}u\allowbreak{}n\allowbreak{}t\allowbreak{}i\allowbreak{}m\allowbreak{}e\allowbreak{}C\allowbreak{}o\allowbreak{}m\allowbreak{}m\allowbreak{}a\allowbreak{}n\allowbreak{}d\allowbreak{}I\allowbreak{}n\allowbreak{}t\allowbreak{}e\allowbreak{}n\allowbreak{}t\allowbreak{}E\allowbreak{}m\allowbreak{}i\allowbreak{}s\allowbreak{}s\allowbreak{}i\allowbreak{}o\allowbreak{}n} values from declared \texttt{C\allowbreak{}o\allowbreak{}m\allowbreak{}m\allowbreak{}a\allowbreak{}n\allowbreak{}d\allowbreak{}I\allowbreak{}n\allowbreak{}t\allowbreak{}e\allowbreak{}n\allowbreak{}t} metadata. Runtime code maps those emissions through explicit \texttt{R\allowbreak{}u\allowbreak{}n\allowbreak{}t\allowbreak{}i\allowbreak{}m\allowbreak{}e\allowbreak{}C\allowbreak{}o\allowbreak{}m\allowbreak{}m\allowbreak{}a\allowbreak{}n\allowbreak{}d\allowbreak{}I\allowbreak{}n\allowbreak{}t\allowbreak{}e\allowbreak{}n\allowbreak{}t\allowbreak{}R\allowbreak{}o\allowbreak{}u\allowbreak{}t\allowbreak{}e} values into the existing \texttt{R\allowbreak{}u\allowbreak{}n\allowbreak{}t\allowbreak{}i\allowbreak{}m\allowbreak{}e\allowbreak{}C\allowbreak{}o\allowbreak{}m\allowbreak{}m\allowbreak{}a\allowbreak{}n\allowbreak{}d\allowbreak{}B\allowbreak{}u\allowbreak{}f\allowbreak{}f\allowbreak{}e\allowbreak{}r}. Business logic still handles the resulting command later through the normal command-buffer path. Emissions copy and deeply freeze inert payload data before routing so caller mutation cannot change the command that eventually reaches business logic.

% Presentation of original arrow notation. See flow-manifest.json.
\begin{center}
\begin{tikzpicture}[book flow node/.style={draw=BookHeading,line width=.65pt,fill=BookPanel,rounded corners=1.5pt,align=center,inner sep=5pt,minimum height=10mm,font=\fontsize{9}{11}\selectfont,execute at begin node={\hyphenpenalty=10000\exhyphenpenalty=10000}},book flow edge/.style={-{Latex[length=2mm]},draw=BookHeading,line width=.7pt}]
\node[book flow node,text width=96mm] (n0) at (0,0.0) {view emits command intent};
\node[book flow node,text width=96mm] (n1) at (0,-1.55) {runtime maps intent to a command};
\node[book flow node,text width=96mm] (n2) at (0,-3.1) {command buffer records the command};
\node[book flow node,text width=96mm] (n3) at (0,-4.65) {business logic applies the command later};
\draw[book flow edge] (n0) -- (n1);
\draw[book flow edge] (n1) -- (n2);
\draw[book flow edge] (n2) -- (n3);
\end{tikzpicture}
\end{center}


Command intents and emissions expose no provider handles, subscription handles, refresh methods, dispatch callbacks, mutable stores, or render hooks.

\subsection{5. User Input Emits Commands}\label{record-132-user-input-emits-commands}

Views communicate user intent through Commands:

\begin{CodeBlock}
command('reader.selectHeading', { headingId });
command('article.refreshRequested', { articleId });
command('nav.openItem', { itemId });
\end{CodeBlock}

Business logic receives those Commands through the existing runtime command path, updates state or provider parameters, and lets providers publish the next snapshot.

Commands are the only sanctioned path from views back toward business logic.

\subsection{6. Binding Facts Lower Across Modes}\label{record-132-binding-facts-lower-across-modes}

Binding status is semantic UI truth.

Ready, loading, empty, stale, and error states must lower across interactive, static, pipe, and accessible modes without parsing rendered text.

Examples:

\begin{itemize}
\tightlist
\item
  loading article body becomes a loading fact and an accessible status line
\item
  stale search results expose stale facts even if the rich view only shows a subtle marker
\item
  failed credential-status provider exposes provider status but never secret values
\item
  empty comments list lowers as an empty-state fact rather than disappearing
\end{itemize}

\subsection{7. Provider Scopes Are Explicit}\label{record-132-provider-scopes-are-explicit}

Provider scopes should be attached by business logic, app shells, tests, or stories:

\begin{CodeBlock}
const app = appShell({
  providers: providerScope([
    provide(articleProvider),
    provide(selectionProvider),
    provide(commandLogProvider),
  ]),
  slots: {
    navigation: navigationBlock(),
    content: readerSurfaceBlock(),
    inspector: inspectorPanelBlock(),
  },
});
\end{CodeBlock}

No import should register a global provider as a side effect.

\section{AppShell Implications}\label{record-132-appshell-implications}

AppShell should not be only a layout helper.

It should be the first proof that:

\begin{itemize}
\tightlist
\item
  named block slots can contain nested provider-bound blocks
\item
  shell-level provider scope can feed navigation, content, inspector, status, and overlay regions
\item
  region bindings are semantic, not physical left/right assumptions
\item
  region Commands bubble as user intent, not direct shell mutation
\item
  active region changes can alter which bindings are active without leaking subscriptions
\end{itemize}

Conceptual AppShell shape:

\begin{CodeBlock}
const docsApp = appShell({
  providers: docsProviderScope,
  slots: {
    navigation: navigationListBlock({
      data: { items: docsNavResource },
      commands: { openItem: 'docs.openArticle' },
    }),
    content: readerSurfaceBlock({
      data: { article: activeArticleResource },
      commands: { selectHeading: 'docs.selectHeading' },
    }),
    inspector: inspectorPanelBlock({
      data: { selection: activeSelectionResource },
    }),
    status: statusBarBlock({
      data: { sync: docsSyncStatusResource },
    }),
  },
});
\end{CodeBlock}

The final API may differ. The architectural point is that slots contain nested blocks that declare their own data requirements, and the active hierarchy resolves those requirements through explicit scopes.

\section{Relationship To Existing Systems}\label{record-132-relationship-to-existing-systems}

\subsection{DX-031 Blocks}\label{record-132-dx-031-blocks}

DX-031 metadata remains useful for discovery, docs, package manifests, and story indexing.

DX-034 adds the missing runtime-authoring layer: declared data requirements, provider scopes, immutable binding frames, and command intent outputs.

\subsection{Runtime Engine}\label{record-132-runtime-engine}

RE-001 through RE-006 already separate state, view stack, retained layouts, input routing, commands, effects, and component interaction.

DX-034 should reuse that spine instead of inventing a second event model.

\subsection{Event Bus}\label{record-132-event-bus}

The event bus is a carrier for runtime messages and command results.

It is not the data-binding source of truth.

A provider may listen to a channel or event stream internally, but the view binding contract still sees immutable snapshots, not the bus.

\subsection{Schema-Bound Blocks}\label{record-132-schema-bound-blocks}

Schema-bound blocks are narrower than provider-bound views.

Schemas validate the shape of data. Providers answer where data comes from, when it updates, what its lifecycle status is, and which active view owns the binding.

The two should compose:

% Presentation of original arrow notation. See flow-manifest.json.
\begin{center}
\begin{tikzpicture}[book flow node/.style={draw=BookHeading,line width=.65pt,fill=BookPanel,rounded corners=1.5pt,align=center,inner sep=5pt,minimum height=10mm,font=\fontsize{9}{11}\selectfont,execute at begin node={\hyphenpenalty=10000\exhyphenpenalty=10000}},book flow edge/.style={-{Latex[length=2mm]},draw=BookHeading,line width=.7pt}]
\node[book flow node,text width=96mm] (n0) at (0,0.0) {provider snapshot};
\node[book flow node,text width=96mm] (n1) at (0,-1.55) {schema adapter validates data};
\node[book flow node,text width=96mm] (n2) at (0,-3.1) {immutable binding frame};
\node[book flow node,text width=96mm] (n3) at (0,-4.65) {block render};
\draw[book flow edge] (n0) -- (n1);
\draw[book flow edge] (n1) -- (n2);
\draw[book flow edge] (n2) -- (n3);
\end{tikzpicture}
\end{center}


\section{Accessibility / Assistive Posture}\label{record-132-accessibility-assistive-posture}

Assistive modes need the binding state, not only the rendered cells.

Every provider-bound view should expose:

\begin{itemize}
\tightlist
\item
  loading state
\item
  stale state
\item
  empty state
\item
  error state
\item
  selected or active entity, when relevant
\item
  available intent Commands, when relevant
\item
  redaction facts for sensitive provider-backed surfaces
\end{itemize}

\section{Localization / Directionality Posture}\label{record-132-localization-directionality-posture}

Providers should deliver data, not localized shell policy.

Localized text can be data when it comes from a catalog provider, but layout direction, fallback copy, and shell-owned labels remain host or i18n runtime responsibility.

Data requirements should avoid physical slot names such as \texttt{l\allowbreak{}e\allowbreak{}f\allowbreak{}t\allowbreak{}N\allowbreak{}a\allowbreak{}v} or \texttt{r\allowbreak{}i\allowbreak{}g\allowbreak{}h\allowbreak{}t\allowbreak{}I\allowbreak{}n\allowbreak{}s\allowbreak{}p\allowbreak{}e\allowbreak{}c\allowbreak{}t\allowbreak{}o\allowbreak{}r}; use logical names such as \texttt{n\allowbreak{}a\allowbreak{}v\allowbreak{}i\allowbreak{}g\allowbreak{}a\allowbreak{}t\allowbreak{}i\allowbreak{}o\allowbreak{}n}, \texttt{c\allowbreak{}o\allowbreak{}n\allowbreak{}t\allowbreak{}e\allowbreak{}n\allowbreak{}t}, \texttt{i\allowbreak{}n\allowbreak{}s\allowbreak{}p\allowbreak{}e\allowbreak{}c\allowbreak{}t\allowbreak{}o\allowbreak{}r}, and \texttt{s\allowbreak{}t\allowbreak{}a\allowbreak{}t\allowbreak{}u\allowbreak{}s}.

\section{Agent Inspectability / Explainability Posture}\label{record-132-agent-inspectability-explainability-posture}

Agents should be able to inspect:

\begin{itemize}
\tightlist
\item
  the active view hierarchy
\item
  each view's data requirements
\item
  provider resolution for each requirement
\item
  current binding snapshot status and version
\item
  Commands each view may emit
\item
  binding invalidation causes
\item
  teardown/suspension state for inactive views
\end{itemize}

Agents should not need to parse rendered terminal output to understand data flow.

\section{Linked Invariants}\label{record-132-linked-invariants}

\begin{itemize}
\tightlist
\item
  \href{https://github.com/flyingrobots/bijou/blob/cbd70612140e7a450d764b3c7be1c6af878b99bc/docs/invariants/runtime-truth-wins.md}{Runtime Truth Wins}
\item
  \href{https://github.com/flyingrobots/bijou/blob/cbd70612140e7a450d764b3c7be1c6af878b99bc/docs/invariants/schemas-live-at-boundaries.md}{Schemas Live At Boundaries}
\item
  \href{https://github.com/flyingrobots/bijou/blob/cbd70612140e7a450d764b3c7be1c6af878b99bc/docs/invariants/commands-change-state-effects-do-not.md}{Commands Change State, Effects Do Not}
\item
  \href{https://github.com/flyingrobots/bijou/blob/cbd70612140e7a450d764b3c7be1c6af878b99bc/docs/invariants/state-machine-and-view-stack-are-distinct.md}{State Machine and View Stack Are Distinct}
\item
  \href{https://github.com/flyingrobots/bijou/blob/cbd70612140e7a450d764b3c7be1c6af878b99bc/docs/invariants/layout-owns-interaction-geometry.md}{Layout Owns Interaction Geometry}
\item
  \href{https://github.com/flyingrobots/bijou/blob/cbd70612140e7a450d764b3c7be1c6af878b99bc/docs/invariants/tests-are-the-spec.md}{Tests Are the Spec}
\end{itemize}

\section{Implementation Outline}\label{record-132-implementation-outline}

\begin{enumerate}
\def\labelenumi{\arabic{enumi}.}
\tightlist
\item
  Done: define the public architecture vocabulary:

  \begin{itemize}
  \tightlist
  \item
    \texttt{D\allowbreak{}a\allowbreak{}t\allowbreak{}a\allowbreak{}R\allowbreak{}e\allowbreak{}q\allowbreak{}u\allowbreak{}i\allowbreak{}r\allowbreak{}e\allowbreak{}m\allowbreak{}e\allowbreak{}n\allowbreak{}t}
  \item
    \texttt{B\allowbreak{}i\allowbreak{}n\allowbreak{}d\allowbreak{}i\allowbreak{}n\allowbreak{}g\allowbreak{}S\allowbreak{}n\allowbreak{}a\allowbreak{}p\allowbreak{}s\allowbreak{}h\allowbreak{}o\allowbreak{}t}
  \item
    \texttt{B\allowbreak{}i\allowbreak{}n\allowbreak{}d\allowbreak{}i\allowbreak{}n\allowbreak{}g\allowbreak{}F\allowbreak{}r\allowbreak{}a\allowbreak{}m\allowbreak{}e}
  \item
    \texttt{C\allowbreak{}o\allowbreak{}m\allowbreak{}m\allowbreak{}a\allowbreak{}n\allowbreak{}d\allowbreak{}I\allowbreak{}n\allowbreak{}t\allowbreak{}e\allowbreak{}n\allowbreak{}t}
  \end{itemize}
\item
  Done: add behavioral tests proving the design remains unidirectional, declarative, and immutable at the primitive layer.
\item
  Done: add pure runtime primitives for immutable binding snapshots and binding frames.
\item
  Done: define provider and provider-scope contracts without global registration.
\item
  Done: add one-scope provider resolution.
\item
  Done: assemble immutable binding frames from resolved snapshots.
\item
  Done: declare view data contracts.
\item
  Done: integrate view data and command contracts with DX-031 block definitions.
\item
  Done: add a structural AppShell composition contract for semantic slots, explicit provider scopes, and nested block data/command introspection.
\item
  Done: formalize the active binding lifecycle before runtime subscriptions or rendered AppShell begin.
\item
  Done: add binding lifecycle primitives as immutable transition-algebra value objects.
\item
  Done: add active binding collection primitives for declared owner and requirement pairs.
\item
  Done: add active-view binding collection over the existing view-stack model.
\item
  Done: add invalidation flow from provider snapshot updates to new immutable binding frames.
\item
  Done: add Command intent dispatch proof through the runtime command buffer.
\item
  Done: prove rendered AppShell with provider-bound navigation, content, inspector, and status blocks.
\item
  Done: add DOGFOOD stories and captures for ready, loading, stale, empty, and error binding states through the standard block story and DOGFOOD preview surface.
\end{enumerate}

\section{Tests To Write First}\label{record-132-tests-to-write-first}

\begin{itemize}
\tightlist
\item
  Behavioral tests proving public binding frames expose immutable data without provider handles.
\item
  Behavioral tests proving binding snapshots are immutable and versioned.
\item
  Behavioral tests proving command intents are metadata, not callbacks.
\item
  Behavioral tests proving provider scopes are explicit local registries without hidden globals.
\item
  Cycle tests proving the active hierarchy owns active binding lifetime, provider subscriptions belong to runtime lifecycle, and binding frames remain immutable render inputs.
\item
  Behavioral tests proving lifecycle records are immutable transition-algebra facts with active, suspended, disposed, and invalidated states.
\item
  Behavioral tests proving active binding collections bind declared requirements to explicit lifecycle owners without rendering, subscribing, or dispatching.
\item
  Runtime tests proving active view-stack layers determine active bindings without rendering, subscribing, or dispatching.
\item
  Runtime tests proving provider scopes resolve nearest-provider wins without hidden globals.
\item
  Runtime tests proving active views create bindings and inactive views dispose or suspend them.
\item
  Runtime tests proving provider updates create new binding frames instead of mutating the current frame.
\item
  Runtime tests proving views can emit Commands but cannot mutate provider data.
\item
  Runtime tests proving Commands flow to business logic and produce later snapshots through the normal update loop.
\item
  Mode-lowering tests proving ready, loading, stale, empty, and error binding states survive in pipe and accessible output.
\item
  DOGFOOD tests proving AppShell nests provider-bound blocks through named logical regions.
\end{itemize}

\section{Acceptance Criteria}\label{record-132-acceptance-criteria}

\begin{itemize}
\tightlist
\item
  Bijou has a documented data-binding architecture before expanding the block API further.
\item
  Data flow is one-way from business logic/providers into immutable view snapshots.
\item
  Views communicate user intent only through Commands.
\item
  Providers are explicit, scoped, and free of import-time global registration.
\item
  Active view hierarchy controls which data bindings are active.
\item
  Active binding collections expose active owner/requirement pairs and lifecycle ownership records before runtime provider integration.
\item
  Provider updates produce new immutable binding frames.
\item
  Command intent emissions enter the runtime command buffer without embedding business logic callbacks in views.
\item
  Binding status is visible to tooling and lower modes.
\item
  Schema-bound blocks compose with provider-bound views but do not replace the provider lifecycle.
\item
  AppShell is designed around nested provider-bound blocks in logical slots.
\end{itemize}

\section{Risks / Unknowns}\label{record-132-risks-unknowns}

\begin{itemize}
\tightlist
\item
  The first API could accidentally expose provider handles to views. That would violate the core rule.
\item
  The provider lifecycle could duplicate existing runtime command/effect buffering if it is not layered carefully.
\item
  If provider scopes are too implicit, debugging active data flow will become difficult.
\item
  If binding frames are only TypeScript shapes, they will not satisfy the repo's runtime-truth standard.
\item
  Schema binding, provider binding, and slot binding can blur together. The architecture should keep them separate:

  \begin{itemize}
  \tightlist
  \item
    schemas validate data shape
  \item
    providers deliver immutable snapshots
  \item
    slots compose views and blocks
  \item
    Commands express user intent
  \end{itemize}
\end{itemize}

\section{Retrospective}\label{record-132-retrospective}

DX-034A landed the primitive contract layer first. The important restraint is that \texttt{B\allowbreak{}i\allowbreak{}n\allowbreak{}d\allowbreak{}i\allowbreak{}n\allowbreak{}g\allowbreak{}F\allowbreak{}r\allowbreak{}a\allowbreak{}m\allowbreak{}e} is not a provider scope, subscription manager, active-view resolver, schema adapter, or AppShell slot model. It is only the immutable data frame views can read during render.

DX-034B then landed explicit provider/scope metadata. The restraint remains the same: \texttt{P\allowbreak{}r\allowbreak{}o\allowbreak{}v\allowbreak{}i\allowbreak{}d\allowbreak{}e\allowbreak{}r\allowbreak{}S\allowbreak{}c\allowbreak{}o\allowbreak{}p\allowbreak{}e} says which providers are locally available, but it still does not perform nearest-scope resolution, subscribe to backing sources, invalidate views, or dispatch Commands.

DX-034C landed structural AppShell composition. It proves semantic slots and nested block introspection, but still does not render AppShell or invent binding lifecycle policy.

DX-034D formalized active binding lifecycle ownership before runtime code. The important restraint is that active hierarchy, provider subscriptions, invalidation, and command dispatch are runtime responsibilities; views still receive immutable frames and emit intent only.

DX-034E landed binding lifecycle primitives as transition algebra. The restraint is that lifecycle records are immutable facts, not managers: they do not subscribe, refresh, dispatch, traverse the active hierarchy, render AppShell, retain caches, or bind schemas.

DX-034F landed active binding collection primitives. The restraint is that collections inspect declarations and produce active lifecycle ownership records; they do not subscribe, refresh, dispatch, render, cache, resolve provider scopes, traverse the real TUI runtime tree, bind schemas, or integrate DOGFOOD.

DX-034G landed active view-stack binding collection in \texttt{@\allowbreak{}f\allowbreak{}l\allowbreak{}y\allowbreak{}i\allowbreak{}n\allowbreak{}g\allowbreak{}r\allowbreak{}o\allowbreak{}b\allowbreak{}o\allowbreak{}t\allowbreak{}s\allowbreak{}/\allowbreak{}b\allowbreak{}i\allowbreak{}j\allowbreak{}o\allowbreak{}u\allowbreak{}-\allowbreak{}t\allowbreak{}u\allowbreak{}i}. The restraint is that the runtime stack only decides which declared binding sources are active; it still does not render, subscribe, refresh, dispatch, resolve provider handles, or retain caches.

DX-034H landed provider-update invalidation and frame assembly as a pure contract helper. The restraint is that snapshot updates create new immutable frames and lifecycle invalidation facts; no provider manager or subscription loop exists yet.

DX-034I landed command intent emission and runtime command-buffer routing. The restraint is that command routes map intent records to command values; business logic still applies commands later through the existing runtime command buffer.

DX-034 is landed for the \texttt{v\allowbreak{}6\allowbreak{}.\allowbreak{}0\allowbreak{}.\allowbreak{}0} release boundary. The shipped release floor is the explicit data-binding contract layer: immutable snapshots and frames, provider scopes, view data contracts, structural AppShell composition, active binding lifecycle facts, active runtime binding collection, provider-update frame assembly, command intent routing, and closeout proof that a provider-bound AppShell can render \texttt{n\allowbreak{}a\allowbreak{}v\allowbreak{}i\allowbreak{}g\allowbreak{}a\allowbreak{}t\allowbreak{}i\allowbreak{}o\allowbreak{}n}, \texttt{c\allowbreak{}o\allowbreak{}n\allowbreak{}t\allowbreak{}e\allowbreak{}n\allowbreak{}t}, \texttt{i\allowbreak{}n\allowbreak{}s\allowbreak{}p\allowbreak{}e\allowbreak{}c\allowbreak{}t\allowbreak{}o\allowbreak{}r}, and \texttt{s\allowbreak{}t\allowbreak{}a\allowbreak{}t\allowbreak{}u\allowbreak{}s} regions from immutable binding frames without exposing provider, subscription, refresh, or dispatch handles to rendered output.

Provider-bound AppShell proof is present in the DX-034 closeout cycle test. The proof intentionally stays pure: provider subscriptions and cache policy remain runtime/provider responsibilities, while the release contract proves the boundary that views consume immutable frames and emit command intents.

DOGFOOD binding-state proof is present through the standard block story states and Blocks preview lowering checks for ready, loading, stale, empty, and error states. Broader production AppShell behavior can now build on the landed contract without keeping issue \#182 open.

